\documentclass[../../../include/open-logic-section]{subfiles}

\begin{document}

\olfileid{sth}{ordinals}{ordtype}
\olsection{Ordinal Minangka Tipe Urutan}

Kanthi Panggantian, mula saiki kita makarya ing
$\ZFminus$, pungkasane kita bisa mbuktekake asil
kang wis kita tuju:

\begin{thm}\ollabel{thmOrdinalRepresentation}
Saben urutan becik isomorfis karo ordinal kang tunggal.
\end{thm}

\begin{proof}
Upamane $\tuple{A, <}$ urutan becik. Miturut
\olref[basic]{ordisoidentity}, urutan iki paling akeh
isomorfis karo siji ordinal. Kanggo kontradiksi,
upamane $\tuple{A, <}$ ora isomorfis karo
\emph{ordinal endi wae}. Pisanan kita bakal
``nggawe $\tuple{A, <}$ sacilik-cilike''.
Rincine: yen ana perangan wiwitan sejati
$\tuple{A_a, <_a}$ kang ora isomorfis karo ordinal,
ana $a \in A$ paling cilik kang nduweni sipat iku;
banjur upamane $B = A_a$ lan
$\mathord{\lessdot} = \mathord{<_a}$.
Yen ora, upamane $B = A$ lan
$\mathord{\lessdot} = \mathord{<}$.

Miturut pilihan iki, saben perangan wiwitan sejati
saka $B$ isomorfis karo sawenehing ordinal kang
tunggal kaya ing ndhuwur. Dadi miturut Panggantian,
himpunan ing ngisor iki ana lan dadi fungsi:
\[
	f = \Setabs{\tuple{\beta, b}}{b \in B\text{ lan }\ordeq{\beta}{\tuple{B_b, \lessdot_b}}}
\]
Kanggo ngrampungake kontradiksi, kita bakal nuduhake
yen $f$ isomorfisme $\alpha \to B$ kanggo sawijining
ordinal $\alpha$.

Cetha yen $\ran{f} = B$. Miturut
\olref[iso]{lemordsegments}, $f$ njaga urutan, yaiku
$\gamma \in \beta$ yen lan mung yen
$f(\gamma) \lessdot f(\beta)$. Kanggo nuduhake yen
$\dom{f}$ ordinal, miturut
\olref[basic]{corordtransitiveord} cukup nuduhake
yen $\dom{f}$ transitif. Pilihen
$\beta \in \dom{f}$, yaiku
$\ordeq{\beta}{\tuple{B_b, \lessdot_b}}$ kanggo
sawijining $b$. Yen $\gamma \in \beta$, mula
$\gamma \in \dom{f}$ miturut
\olref[iso]{wellordinitialsegment}; kanthi umum,
$\beta \subseteq \dom{f}$.
\end{proof}

Asil iki ngidini definisi kang wis dikarepake wiwit
\olref[vn]{sec}:

\begin{defn}
Yen $\tuple{A, <}$ urutan becik, tipe urutane,
$\ordtype{A, <}$, yaiku ordinal tunggal $\alpha$
kang nyukupi $\ordeq{\tuple{A, < }}{\alpha}$.
\end{defn}

Kajaba iku, definisi iki ngasilake rong prinsip:

\begin{cor}\ollabel{ordtypesworklikeyouwant}
Yen $\tuple{A, <}$ lan $\tuple{B, \lessdot}$
urutan becik:
\begin{align*}
	\ordtype{A, <} = \ordtype{B, \lessdot}&\text{ yen lan mung yen }\ordeq{\tuple{A, <}}{\tuple{B, \lessdot}}\\
	\ordtype{A, <} \in \ordtype{B, \lessdot}&\text{ yen lan mung yen }\ordeq{\tuple{A, <}}{\tuple{B_b, \lessdot_b}}\text{ kanggo sawenehing }b \in B
\end{align*}
\end{cor}

\begin{proof}
Persamaan kapisan ndherek saka
\olref[basic]{ordisoidentity}. Kanggo mbuktekake
pratelan kapindho, upamane $\ordtype{A, <} = \alpha$
lan $\ordtype{B, \lessdot} = \beta$, lan upamane
$f \colon \beta \to B$ isomorfisme kita. Mula:
\begin{align*}
	\alpha \in \beta&\Longrightarrow\funrestrictionto{f}{\alpha} \colon \alpha \to B_{f(\alpha)}\text{ iku isomorfisme}\\
	&\Longrightarrow\ordeq{\tuple{A, <}}{\tuple{B_{f(\alpha)}, \lessdot_{f(\alpha)}}}\\
	&\Longrightarrow\ordeq{\tuple{A, <}}{\tuple{B_b, \lessdot_b}}\text{ kanggo sawenehing $b \in B$}
\end{align*}
miturut \olref[iso]{ordisounique},
\olref[iso]{wellordinitialsegment}, lan
\olref[basic]{ordissetofsmallerord}. Kosok baline,
yen urutan kapisan isomorfis karo sawijining perangan
wiwitan sejati saka urutan kapindho, titik pungkasan
perangan iku nduweni prabayangan ing ordinal sumber
miturut isomorfisme kasebut. Perangan mau isomorfis
karo prabayangan iki minangka ordinal; mula tipe
urutan kapisan padha karo prabayangan kasebut miturut
\olref[basic]{ordisoidentity}. Dadi tipe urutan kapisan
iku anggota ordinal sumber mau.
\end{proof}

\end{document}
